\section{Post-training III：RLVR 与可验证奖励}

\subsection{为什么 RLVR 成为推理提升的关键转折}
讲者将 RLVR（Reinforcement Learning with Verifiable Rewards）定义为后训练中的关键拐点：当任务可以自动验证正确与否时，模型可以从大量自生成样本中持续改进，而不再完全依赖昂贵的人类偏好标注。

\begin{importantbox}{RLVR 的最小闭环}
\begin{enumerate}
    \item 生成候选答案；
    \item 用规则或程序验证器给出可验证奖励；
    \item 用 RL 更新策略，提高高奖励样本概率。
\end{enumerate}
\end{importantbox}

\begin{figure}[H]
\centering
\includegraphics[width=0.9\textwidth]{frame-06-rlvr.jpg}
\caption{视频约 01:03:20：从神经 reward model 转向可验证奖励，降低标注依赖}
\end{figure}

\subsection{神经奖励模型与规则奖励的比较}
\begin{knowledgebox}{为何讲者强调 ``可验证''}
神经奖励模型灵活但存在偏差与不可解释性；规则奖励更窄但更可信。对数学、代码、符号推理等任务，规则奖励能显著降低 reward hacking 风险。
\end{knowledgebox}

\begin{table}[H]
\centering
\renewcommand{\arraystretch}{1.2}
\begin{tabular}{p{2.8cm}p{3.6cm}p{3.8cm}p{2.8cm}}
\toprule
\textbf{奖励类型} & \textbf{优点} & \textbf{缺点} & \textbf{典型任务} \\
\midrule
神经奖励模型 & 表达能力强，可覆盖复杂主观目标 & 易偏置、可解释性弱、成本高 & 对话偏好、风格一致性 \\
规则/可验证奖励 & 稳定、可审计、可自动化规模训练 & 覆盖范围受限，设计成本高 & 数学、代码、形式化约束 \\
\bottomrule
\end{tabular}
\caption{两类奖励机制对比}
\end{table}

\begin{warningbox}{RLVR 的边界条件}
RLVR 并不自动提升所有任务。若任务缺少可靠验证器，或者验证器与真实目标错位，训练会把模型推向 ``高分低质'' 的错误方向。
\end{warningbox}

\subsection{从 DeepSeek-R1 经验到开放社区可复现实践}
讲者引用了 R1 风格路线的启发：当验证器便宜、样本可规模生成时，RL 能迅速放大推理能力。对开放社区而言，关键在于把验证器、训练脚本和评估流程一起开源，否则很难验证结论可迁移性。

\subsection{本章小结}
RLVR 的真正价值在于把 ``推理提升'' 从高度人工依赖转为 ``可程序化迭代''。前提是验证器质量可靠且与真实目标一致。


